\chapter{THAM SỐ, KÝ HIỆU VÀ CẤU HÌNH THỰC NGHIỆM}
\label{app:config}

\section{Danh mục ký hiệu toán học}
\begin{longtable}{p{3.0cm}p{11.0cm}}
\caption{Danh mục các ký hiệu chính}\label{tab:symbols}\\
\toprule
\textbf{Ký hiệu} & \textbf{Ý nghĩa}\\
\midrule
\endfirsthead
\toprule
\textbf{Ký hiệu} & \textbf{Ý nghĩa}\\
\midrule
\endhead
$K$ & Số người dùng\\
$N$ & Số phần tử STAR--RIS\\
$h_{d,k}$ & Kênh trực tiếp BS--user $k$\\
$\bm g$ & Kênh BS--STAR--RIS\\
$\bm h_{r,k}$ & Kênh STAR--RIS--user $k$\\
$u_k$ & Chỉ báo phía user: 0 phản xạ, 1 truyền\\
$\beta_n^T,\beta_n^R$ & Tỷ lệ công suất truyền và phản xạ của phần tử $n$\\
$\theta_n^T,\theta_n^R$ & Pha truyền và phản xạ\\
$\phi_n^T,\phi_n^R$ & Hệ số phức truyền và phản xạ\\
$h_k^{\mathrm{eff}}$ & Kênh hiệu dụng của user $k$\\
$p_c,p_k$ & Công suất luồng chung và luồng riêng $k$\\
$\eta_k$ & Tỷ lệ common rate cấp cho user $k$\\
$R_c$ & Common rate mà mọi user có thể giải mã\\
$R_{p,k}$ & Private rate của user $k$\\
$R_k$ & Tổng tốc độ user $k$\\
$R_{\min}$ & Ngưỡng QoS tối thiểu\\
$V,V_2$ & Tổng deficit và deficit bình phương\\
$\lambda_t$ & Dual multiplier cho QoS\\
$d_s,d_a$ & Số chiều state và action\\
$\gamma$ & Discount factor\\
$\tau$ & Hệ số soft target update\\
\bottomrule
\end{longtable}

\section{Cấu hình theo kích thước STAR--RIS}
Các tệp cấu hình phiên bản cuối nằm trong \texttt{configs/v3/}. Mọi cấu hình giữ $K=4$ và thay đổi $N$. Bảng~\ref{tab:dimension-by-n} trình bày số chiều theo công thức source.

\begin{table}[H]
\caption{Số chiều observation và action theo $N$}
\label{tab:dimension-by-n}
\centering
\begin{tabular}{rrrr}
\toprule
$N$ & $d_s=2(K+N+KN)+K$ & $d_a=2K+1+3N$ & Tổng biến beta/pha\\
\midrule
16 & 140 & 57 & 48\\
32 & 268 & 105 & 96\\
64 & 524 & 201 & 192\\
96 & 780 & 297 & 288\\
128 & 1036 & 393 & 384\\
\bottomrule
\end{tabular}
\tablesource{Tác giả tính từ implementation (2026)}
\end{table}

\section{Tệp cấu hình mẫu}
Đoạn dưới là cấu hình $N=32$ được rút gọn để người đọc đối chiếu với Chương~\ref{chap:method}.
\begin{lstlisting}[style=thesispython,caption={Cấu hình TD3 constrained cho $N=32$},label={lst:config-n32}]
n_ris: 32
n_users: 4
p_max: 1.0
noise_power: 0.001
qos_min: 0.5
episode_length: 32
gamma: 0.99
tau: 0.005
hidden_dim: 256
batch_size: 256
replay_size: 200000
warmup_steps: 2000
train_steps: 100000
validation_interval: 5000
validation_scenarios: 1000
exploration_noise: 0.12
exploration_noise_final: 0.03
observation_normalization: blockwise_v2
action_parameterization: physical_v3
qos_penalty_linear: 8.0
qos_penalty_quadratic: 8.0
qos_dual_enabled: true
td3_actor_lr: 0.0001
td3_critic_lr: 0.0003
td3_policy_delay: 2
td3_target_noise: 0.15
td3_noise_clip: 0.30
td3_gradient_clip_norm: 10.0
td3_noise_reference_dim: 64
td3_critic_loss: huber
td3_layer_norm: true
\end{lstlisting}

\section{Quy tắc số và đơn vị}
Trong nội dung tiếng Việt, số thập phân được trình bày bằng dấu phẩy và nhóm ba chữ số bằng dấu chấm theo hướng dẫn của trường. Trong tệp cấu hình và mã nguồn, dấu chấm vẫn được dùng vì là cú pháp số thực của Python/YAML. Đơn vị tốc độ phổ là bit/s/Hz; latency dùng ms hoặc $\mu$s tùy bảng raw.
